%%--------------------------------------------------------------------
%1--------------------------------------------------------------------
\begin{glyph}{Painting (画)}{painting}\label{glyph:painting}
Painting.\\\hline
A secondary meaning has to do with ``漢字'' strokes (think of how they are painted in calligraphy), and things being \textit{drawn up}, from dividing lines to plans.
\end{glyph}
\gindex{Painting}{画}\strindex{8}{画}

\begin{comps}
\comprtwocone&由 + 一　+ 凵\\\hline
\comprthreecone&Cause (\ref{glyph:cause}) + One + Wide-opened mouth\\\hline
\end{comps}

\begin{remglyph}
\hooglig{Paintings} made by that \remcom{one} caused a \remcom{wide-opened mouth}.\\\hline
\end{remglyph}

\begin{prons}
\pronsrtwocone&\textbf{ガ、カク (GA, KAKU)}&---\\\hline
\rye{3}{ガ is used when this 漢字 refers to painting.\\\hline
カク is used it refers to such things as 漢字 strokes and plans.}
\pronsrthreecone&---&---\\\hline
\end{prons}

\begin{remprons}
\hooglig{Painting} \comon{kakuro} with \comon{gushes} of paint.\\\hline
{\warn}\textit{Kakuro} is a crossword puzzle game.\\\hline
\end{remprons}

%{\middel}
%&\mbox{()}\\\hline
\begin{somme}
画家&painting + house = \textit{painter; artist}&ガ{\middel}カ \mbox{(GA{\middel}KA)}\\\hline
計画&measure + painting = \textit{plan; scheme}&ケイ{\middel}カク \mbox{(KEI{\middel}KAKU)}\\\hline
画面&painting + mask = \textit{scene; picture; photo}&ガ{\middel}メン \mbox{(GA{\middel}MEN)}\\\hline
\notyet{作画}&make + painting = \textit{drawing pictures; taking photos}&サク{\middel}ガ \mbox{(SAKU{\middel}GA)}\\\hline
\notyet{映画}&reflect + painting = \textit{movie}&エイ{\middel}ガ \mbox{(EI{\middel}GA)}\\\hline
\end{somme}

\begin{decomp}{4}
映画&に&行く&の？\\\hline
エイガ&に&いく&の\\\hline
movie&に&go&の\\\hline
\multicolumn{4}{|r|}{\textit{You're going to a movie?}}\\\hline
\end{decomp}

\end{multicols}
%%--------------------------------------------------------------------
%2--------------------------------------------------------------------
\begin{multicols}{2}
\begin{glyph}{Possible (可)}{possible}\label{glyph:possible}
Possible.\\\hline
It is important to learn the proper stroke order of this 漢字 so as not to confuse it with similar-looking characters we will meet in the future.
\end{glyph}
\gindex{Possible}{可}\strindex{5}{可}

\begin{comps}
\comprtwocone&口 + 丁\\\hline
\comprthreecone&Mouth + Block (\ref{glyph:block})\\\hline
\end{comps}

\begin{remglyph}
\hooglig{Possibly} a \remcom{vampire} in this \remcom{block}.\\\hline
\end{remglyph}

\begin{prons}
\pronsrtwocone&\textbf{カ (KA)}&---\\\hline
\pronsrthreecone&---&---\\\hline
\end{prons}

\begin{remprons}
\hooglig{Possibly} \comon{cussing}.\\\hline
\end{remprons}

%{\middel}
%&()\\\hline
\begin{somme}
可愛い&possible + love = \textit{cute}&カ{\middel}ワイ{\middel}い (\mbox{KA{\middel}WAI\middel{i}})\\\hline
不可欠&not + possible + lack = \textit{indispensable}&フ{\middel}カ{\middel}ケツ \mbox{(FU{\middel}KA{\middel}KETSU)}\\\hline
\notyet{可動}&possible + move = \textit{movable; mobile}&カ{\middel}ドウ (KA{\middel}D\={O})\\\hline
\end{somme}

\begin{decomp}{4}
この&赤ちゃん&とっても&可愛い！\\\hline
この&あかちゃん&とっても&カワイい\\\hline
this&baby&very&cute\\\hline
\multicolumn{4}{|r|}{\textit{This baby is too cute!}}\\\hline
\end{decomp}

\end{multicols}
\pagebreak
%%--------------------------------------------------------------------
%3--------------------------------------------------------------------
\vspace*{\fill}
\begin{multicols}{2}
\begin{glyph}{God (神)}{god}\label{glyph:god}
God.
\end{glyph}
\gindex{God}{神}\strindex{9}{神}

\begin{comps}
\comprtwocone&ネ + 申\\\hline
\comprthreecone&Altar/Display + Declare\\\hline
\end{comps}

\begin{remglyph}
\hooglig{God's} \remcom{altar} was \remcom{declared}.\\\hline
\end{remglyph}

\begin{prons}
\pronsrtwocone&\textbf{シン、ジン (SHIN, JIN)}&\textbf{かみ (kami)}\\\hline
\rye{3}{シン is by far the most common of these readings, although the others do show up in some important words.}
\pronsrthreecone&---&かん、こう (kan, k\={o})\\\hline
\end{prons}

\begin{remprons}
\hooglig{God's} \comon{sheen} \comon{jinxed} the \comkun{kamikaze} with \ucomkun{Kung Fu} to its \ucomkun{core}.\\\hline
\end{remprons}

%{\middel}
%&()\\\hline
\begin{somme}
神&\textit{God}&かみ \mbox{(kami)}\\\hline
神話&God + speak = \textit{myth; mythology}&シン{\middel}ワ \mbox{(SHIN{\middel}WA)}\\\hline
\notyet{精神}&refined + God = \textit{mind; soul; heart}&セイ{\middel}シン \mbox{(SEI{\middel}SHIN)}\\\hline
\notyet{神社}&God + company = \textit{Shinto shrine}&ジン{\middel}ジャ \mbox{(JIN{\middel}JA)}\\\hline
\notyet{神経}&God + via = \textit{nerves; sensitivity}&シン{\middel}ケイ \mbox{(SHIN{\middel}KEI)}\\\hline
\end{somme}

\begin{decomp}{5}
私&は&神&を&信じます。\\\hline
わたし&は&かみ&を&シンじます\\\hline
I&は&God&を&believe\\\hline
\multicolumn{5}{|r|}{\textit{I believe in God.}}\\\hline
\end{decomp}

\end{multicols}
\vhrulefill{2pt}
%%--------------------------------------------------------------------
%4--------------------------------------------------------------------
\begin{multicols}{2}
\begin{glyph}{Perish (亡)}{perish}\label{glyph:perish}
Perish.  Die.
\end{glyph}
\gindex{Perish}{亡}\strindex{3}{亡}

\begin{comps}
\comprtwocone&亡\\\hline
\comprthreecone&Conceal/Dead\\\hline
\end{comps}

\begin{remglyph}
Part of the 漢字 radicals (\#{23}).\\\hline
\end{remglyph}

\begin{prons}
\pronsrtwocone&\textbf{ボウ (B\={O})}&\textbf{な (na)}\\\hline
\pronsrthreecone&モウ (M\={O})&ほろ (horo)\\\hline
\end{prons}

\begin{remprons}
The \hooglig{perishing} \ucomon{maw} is used in \ucomkun{horology}, but it can be a \comkun{numbing}, \comon{boring} experience.\\\hline
{\warn}\textit{Horology} is the art of making clocks and watches.\\\hline
\end{remprons}

%{\middel}
%&()\\\hline
\begin{somme}
亡くなる&\textit{die; pass away}&な{\middel}くなる \mbox{(na{\middel}kunaru)}\\\hline
死亡&death + perish = \textit{death}&シ{\middel}ボウ \mbox{(SHI{\middel}B\={O})}\\\hline
亡国&perish + country = \textit{ruined country}&ボウ{\middel}コク \mbox{(B\={O}{\middel}KOKU)}\\\hline
亡き後&perish + after = \textit{after one's death}&な{\middel}きあと　\mbox{(na{\middel}ki ato)}\\\hline
\notyet{未亡人}&not yet + perish + person = \textit{widow}&ミ{\middel}ボウ{\middel}ジン \mbox{(MI{\middel}B\={O}{\middel}JIN)}\\\hline
\end{somme}
\end{multicols}

\begin{decomp}{7}
多く&の&人&が&事故&で&死亡した。\\\hline
おおく&の&ひと&が&ジコ&で&シボウした。\\\hline
many&の&person&が&accident&で&pass away\\\hline
\multicolumn{7}{|r|}{\textit{Many people died in the accident.}}\\\hline
\end{decomp}
\vspace*{\fill}
\pagebreak
%%--------------------------------------------------------------------
%5--------------------------------------------------------------------
\vspace*{\fill}
\begin{multicols}{2}
\begin{glyph}{Occupy (占)}{occupy}\label{glyph:occupy}
Occupy.\\\hline
Fortune Telling.
\end{glyph}
\gindex{Occupy}{占}\strindex{5}{占}

\begin{comps}
\comprtwocone&卜 + 口\\\hline
\comprthreecone&Divination + Mouth\\\hline
\end{comps}

\begin{remglyph}
\hooglig{Occupied} with a \remcom{divining} \remcom{mouth} about \remcom{vampires}.\\\hline
\end{remglyph}

\begin{prons}
\pronsrtwocone&\textbf{セン (SEN)}&\textbf{うらな (urana)}\\\hline
\pronsrthreecone&---&し (shi)\\\hline
\end{prons}

\begin{remprons}
\hooglig{Occupied} with \comkun{your analysis}, the \comon{central} team had a \ucomkun{shift} in mind.\\\hline
\end{remprons}

%{\middel}
%&()\\\hline
\begin{somme}
占う&\textit{tell fortunes}&うらな{\middel}う \mbox{(urana{\middel}u)}\\\hline
占有&occupy + have = \textit{occupy}&セン{\middel}ユウ \mbox{(SEN{\middel}Y\={U})}\\\hline
\notyet{占星術}&occupy + star + technique = \textit{astrology}&セン{\middel}セイ{\middel}ジュツ \mbox{(SEN{\middel}SEI{\middel}JUTSU)}\\\hline
夢占い{\notcov}&dream + occupy = \textit{oneiromancy; dream fortune-telling}&ゆ{\middel}め　うらな{\middel}い \mbox{(yu{\middel}me urana{\middel}i)}\\\hline
\end{somme}

\begin{decomp}{5}
占い師&に&手相&を&見せました。\\\hline
うらなシ&に&てソウ&を&みせました\\\hline
fortune teller&に&palm reading&を&to show\\\hline
\multicolumn{5}{|r|}{\textit{A fortune-teller read my hand.}}\\\hline
\end{decomp}

\end{multicols}
\vhrulefill{2pt}
%%--------------------------------------------------------------------
%6--------------------------------------------------------------------
\begin{multicols}{2}
\begin{glyph}{Society (世)}{society}\label{glyph:society}
Society.\\\hline
It finds expression in terms such as \textit{the age}, \textit{the times}, and \textit{the era}.
\end{glyph}
\gindex{Society}{世}\strindex{5}{世}

\begin{comps}
\comprtwocone&世\\\hline
\comprthreecone&Society/Generation/World\\\hline
\end{comps}

\begin{remglyph}
Uncommon 漢字 radical.\\\hline
\end{remglyph}

\begin{prons}
\pronsrtwocone&\textbf{セ、セイ (SE, SEI)}&---\\\hline
\pronsrthreecone&---&よ (yo)\\\hline
\end{prons}

\begin{remprons}
\hooglig{Society's} \ucomkun{yomp} of the \comon{septic} \comon{sail}.\\\hline
{\warn}\textit{Yomp} is a march over difficult terrain with heavy equipment.\\\hline
\end{remprons}

%{\middel}
%&()\\\hline
\begin{somme}
世話&society + speak = \textit{help; service}&セ{\middel}ワ \mbox{(SE{\middel}WA)}\\\hline
世間&society + interval = \textit{society; life}&セ{\middel}ケン \mbox{(SE{\middel}KEN)}\\\hline
中世&middle + society = \textit{the Middle Ages}&チュウ{\middel}セイ \mbox{(CH\={U}{\middel}SEI)}\\\hline
\notyet{世界}&society + world = \textit{the world}&セ{\middel}カイ \mbox{(SE{\middel}KAI)}\\\hline
\notyet{世代}&society + substitute = \textit{generation}&セ{\middel}ダイ \mbox{(SE{\middel}DAI)}\\\hline
\end{somme}

\begin{decomp}{5}
息子&を&世間&に&出した。\\\hline
むすこ&を&セケン&に&だした\\\hline
son&を&world&に&to put out\\\hline
\multicolumn{5}{|r|}{\textit{He launched his son in the world.}}\\\hline
\end{decomp}

\end{multicols}
\vspace*{\fill}
\pagebreak
%%--------------------------------------------------------------------
%7--------------------------------------------------------------------
\vspace*{\fill}
\begin{multicols}{2}
\begin{glyph}{Fixed (定)}{fixed}\label{glyph:fixed}
Fixed in time or place.\\\hline
Regular.  Definite.
\end{glyph}
\gindex{Fixed}{定}\strindex{8}{定}

\begin{comps}
\comprtwocone&宀 + 疋\\\hline
\comprthreecone&Roof + Roll/Bolt of Cloth\\\hline
\end{comps}

\begin{remglyph}
\hooglig{Fixing} the \remcom{roof} \remcom{bolts}.\\\hline
\end{remglyph}

\begin{prons}
\pronsrtwocone&\textbf{テイ (TEI)}&\textbf{さだ (sada)}\\\hline
\pronsrthreecone&ジョウ (J\={O})&---\\\hline
\end{prons}

\begin{remprons}
\hooglig{Fixing} \comon{tape} \comkun{suddenly} to his \ucomon{jaw}.\\\hline
\end{remprons}

%{\middel}
%&()\\\hline
\begin{somme}
定まる (intr.)&\textit{to be fixed/determined}&さだ{\middel}まる \mbox{(sada{\middel}maru)}\\\hline
定める (tr.)&\textit{to decide/determine (sth.)}&さだ{\middel}める \mbox{(sada{\middel}meru)}\\\hline
安定&ease + fixed = \textit{stability}&アン{\middel}テイ \mbox{(AN{\middel}TEI)}\\\hline
\notyet{定食}&fixed + eat = \textit{set meal}&テイ{\middel}ショク \mbox{(TEI{\middel}SHOKU)}\\\hline
\notyet{指定}&finger + fixed = \textit{designation}&シ{\middel}テイ \mbox{(SHI{\middel}TEI)}\\\hline
\end{somme}

\begin{decomp}{2}
本日&定休日。\\\hline
ホンジツ&テイキュウび\\\hline
Today&regular holiday\\\hline
\multicolumn{2}{|r|}{\textit{Today is our regular closing day.}}\\\hline
\end{decomp}

\end{multicols}
\vhrulefill{2pt}
%%--------------------------------------------------------------------
%8--------------------------------------------------------------------
\begin{multicols}{2}
\begin{glyph}{Paper (紙)}{paper}\label{glyph:paper}
Paper.
\end{glyph}
\gindex{Paper}{紙}\strindex{10}{紙}

\begin{comps}
\comprtwocone&糸 + 氏\\\hline
\comprthreecone&Thread + Surname\\\hline
\end{comps}

\begin{remglyph}
\hooglig{Paper} \remcom{threadings} are unique to different \remcom{clan}.\\\hline
\end{remglyph}

\begin{prons}
\pronsrtwocone&\textbf{シ (SHI)}&\textbf{かみ (kami)}\\\hline
\pronsrthreecone&---&---\\\hline
\end{prons}

\begin{remprons}
The \hooglig{paper} \comkun{kami} \comon{shielded} us.\\\hline
\end{remprons}

%{\middel}
%&()\\\hline
\begin{somme}
紙&\textit{paper}&かみ \mbox{(kami)}\\\hline
手紙&hand + paper = \textit{letter}&て{\middel}がみ \mbox{(te{\middel}gami)}\\\hline
紙面&paper + mask = \textit{(page) space}&シ{\middel}メン \mbox{(SHI{\middel}MEN)}\\\hline
\notyet{和紙}&harmony (Japan) + paper = \textit{Japanese paper}&ワ{\middel}シ \mbox{(WA{\middel}SHI)}\\\hline
\notyet{用紙}&utilize + paper = \textit{sheets of paper}&ヨウ{\middel}シ \mbox{(Y\={O}{\middel}SHI)}\\\hline
\end{somme}

\begin{decomp}{3}
手紙&を&下さい。\\\hline
てがみ&を&ください\\\hline
letter&を&please\\\hline
\multicolumn{3}{|r|}{\textit{Let me hear from you.}}\\\hline
\end{decomp}

\begin{decomp}{5}
此の&用紙&に&記入し&なさい。\\\hline
この&ヨウシ&に&キニュウし&なさい\\\hline
this&form&に&fill in&to do\\\hline
\multicolumn{5}{|r|}{\textit{Fill in this form.}}\\\hline
\end{decomp}

\end{multicols}
\vspace*{\fill}
\pagebreak
%%--------------------------------------------------------------------
%9--------------------------------------------------------------------
\vspace*{\fill}
\begin{multicols}{2}
\begin{glyph}{Appearance (采)}{appearance}\label{glyph:appearance}
Appearance.\\\hline
The look of something.
\end{glyph}
\gindex{Appearance}{采}\strindex{8}{采}

\begin{comps}
\comprtwocone&爫 + 木\\\hline
\comprthreecone&Claw + Tree\\\hline
\end{comps}

\begin{remglyph}
\hooglig{Appearance} of a \remcom{clawed} \remcom{tree}.\\\hline
\end{remglyph}

\begin{prons}
\pronsrtwocone&\textbf{サイ (SAI)}&---\\\hline
\pronsrthreecone&---&---\\\hline
\end{prons}

\begin{remprons}
\hooglig{Appear} \comon{psyched} up.\\\hline
\end{remprons}

%{\middel}
%&()\\\hline
\begin{somme}
采&\textit{appearance}&サイ \mbox{(SAI)}\\\hline
神采&God + appearance = \textit{exceptional appearance}&シン{\middel}サイ \mbox{(SHIN{\middel}SAI)}\\\hline
\notyet{風采}&wind + appearance = \textit{appearance; bearing}&フウ{\middel}サイ \mbox{(F\={U}{\middel}SAI)}\\\hline
\end{somme}
\end{multicols}

\begin{decomp}{7}
ヂック&は&風采&が&母親&に&似ている。\\\hline
ヂック&は&フウサイ&が&ははおや&に&にている\\\hline
Dick&は&appearance&が&mother&に&to resemble\\\hline
\multicolumn{7}{|r|}{\textit{Dick takes after his mother in appearance.}}\\\hline
\end{decomp}
\vhrulefill{2pt}
%%--------------------------------------------------------------------
%10-------------------------------------------------------------------
\begin{multicols}{2}
\begin{glyph}{Army (軍)}{army}\label{glyph:army}
Army.  Military.
\end{glyph}
\gindex{Army}{軍}\strindex{9}{軍}

\begin{comps}
\comprtwocone&冖 + 車\\\hline
\comprthreecone&Cover/Crown + Car\\\hline
\end{comps}

\begin{remglyph}
\hooglig{The army} took \remcom{cover} behind the \remcom{car}.\\\hline
\end{remglyph}

\begin{prons}
\pronsrtwocone&\textbf{グン (GUN)}&---\\\hline
\pronsrthreecone&---&---\\\hline
\end{prons}

\begin{remprons}
\hooglig{Army} of \comon{goons}.\\\hline
{\warn}\textit{Goon} refers to a fool.\\\hline
\end{remprons}

%{\middel}
%&()\\\hline
\begin{somme}
軍人&army + person = \textit{military person; soldier}&グン{\middel}ジン \mbox{(GUN{\middel}JIN)}\\\hline
海軍&sea + army = \textit{navy}&カイ{\middel}グン \mbox{(KAI{\middel}GUN)}\\\hline
空軍&empty (sky) + army = \textit{air force}&クウ{\middel}グン \mbox{(K\={U}{\middel}GUN)}\\\hline
軍事&army + matter = \textit{military affairs}&グン{\middel}ジ　\mbox{(GUN{\middel}JI)}\\\hline
軍団&army + group = \textit{army corps}&グン{\middel}ダン \mbox{(GUN{\middel}DAN)}\\\hline
\notyet{軍用}&army + utilize = \textit{for military use}&グン{\middel}ヨウ \mbox{(GUN{\middel}Y\={O})}\\\hline
\end{somme}
\end{multicols}

\begin{decomp}{7}
彼&の&祖父&は&高級&軍人&だった。\\\hline
かれ&の&ソフ&は&コウキュウ&グンジン&だった\\\hline
he&の&grandfather&は&high rank&soldier&だった\\\hline
\multicolumn{7}{|r|}{\textit{His grandfather was a soldier of high degree.}}\\\hline
\end{decomp}
\vspace*{\fill}
\pagebreak
%%--------------------------------------------------------------------
%11-------------------------------------------------------------------
\vspace*{\fill}
\begin{multicols}{2}
\begin{glyph}{Elder Sister (姉)}{elder_sister}\label{glyph:elder_sister}
Elder sister.
\end{glyph}
\gindex{Elder Sister}{姉}\strindex{8}{姉}

\begin{comps}
\comprtwocone&女 + 市\\\hline
\comprthreecone&Woman + City (\ref{glyph:city})\\\hline
\end{comps}

\begin{remglyph}
\hooglig{Sisters} of \remcom{women} from that \remcom{city}.\\\hline
\end{remglyph}

\begin{prons}
\pronsrtwocone&\textbf{シ (SHI)}&\textbf{あね (ane)}\\\hline
\pronsrthreecone&---&---\\\hline
\end{prons}

\begin{remprons}
\hooglig{Sisterly} \comkun{anemone} \comon{shielded} us.\\\hline
{\warn}\textit{Anemone} are marine, predatory animals.\\\hline
\end{remprons}

%{\middel}
%&()\\\hline
\begin{somme}
姉&\textit{elder sister}&あね \mbox{(ane)}\\\hline
\notyet{姉妹}&elder sister + younger sister = \textit{sister (city/company, etc.)}&シ{\middel}マイ \mbox{(SHI{\middel}MAI)}\\\hline
\notyet{姉妹都市}&sister + metropolis + city = \textit{sister city}&シ{\middel}マイ{\middel}ト{\middel}シ \mbox{(SHI{\middel}MAI{\middel}TO{\middel}SHI)}\\\hline
\notyet{実姉}&truth + sister = \textit{biological sister}&ジッ{\middel}シ \mbox{(JIS{\middel}SHI)}\\\hline
\notyet{長姉}&long + sister = \textit{eldest sister}&チョウ{\middel}シ \mbox{(CH\={O}{\middel}SHI)}\\\hline
\end{somme}

\begin{irrr}
お姉さん&\textit{elder sister}&お{\middel}ねえ{\middel}さん \mbox{(o{\middel}n\={e}{\middel}san)}\\\hline
\rye{3}{As we saw with a similar word for 兄, this is a more polite form of address.}
\end{irrr}

\begin{decomp}{4}
姉妹&が&います&か。\\\hline
シマイ&が&います&か\\\hline
sisters&が&to be&か\\\hline
\multicolumn{4}{|r|}{\textit{Do you have any sisters?}}\\\hline
\end{decomp}

\end{multicols}
\vhrulefill{2pt}
%%--------------------------------------------------------------------
%12-------------------------------------------------------------------
\begin{multicols}{2}
\begin{glyph}{Sheep (羊)}{sheep}\label{glyph:sheep}
Sheep.
\end{glyph}
\gindex{Sheep}{羊}\strindex{6}{羊}

\begin{comps}
\comprtwocone&羊\\\hline
\comprthreecone&Sheep\\\hline
\end{comps}

\begin{remglyph}
Part of the 漢字 radicals (\#{123}).\\\hline
\end{remglyph}

\begin{prons}
\pronsrtwocone&---&\textbf{ひつじ (hitsuji)}\\\hline
\pronsrthreecone&ヨウ (Y\={O})&---\\\hline
\end{prons}

\begin{remprons}
The \hooglig{sheep} \comkun{hits Jeeps} while \ucomon{yawning}.\\\hline
{\warn}\textit{Jeep} is a brand of American automobile.\\\hline
\end{remprons}

%{\middel}
%&()\\\hline
\begin{somme}
羊&\textit{sheep}&ひつじ \mbox{(hitsuji)}\\\hline
子羊&child + sheep = \textit{lamb}&こ{\middel}ひつじ \mbox{(ko{\middel}hitsuji)}\\\hline
羊毛&sheep + fur = \textit{wool}&ヨウ{\middel}モウ \mbox{(Y\={O}{\middel}M\={O})}\\\hline
羊肉&sheep + meat = \textit{lamb meat}&ヨウ{\middel}ニク \mbox{(Y\={O}{\middel}NIKU)}\\\hline
\end{somme}

\begin{decomp}{5}
子羊&は&狼&に&殺された。\\\hline
こひつじ&は&おおかみ&に&ころされた\\\hline
lamb&は&wolf&に&to kill\\\hline
\multicolumn{5}{|r|}{\textit{The lamb was killed by the wolf.}}\\\hline
\end{decomp}

\end{multicols}
\vspace*{\fill}
\pagebreak
%%--------------------------------------------------------------------
%13-------------------------------------------------------------------
\begin{multicols}{2}
\begin{glyph}{Intervene (介)}{intervene}\label{glyph:intervene}
Intervene.  Mediate.
\end{glyph}
\gindex{Intervene}{介}\strindex{4}{介}

\begin{comps}
\comprtwocone&人 + 丿 + 丨\\\hline
\comprthreecone&Person + Slash/NO + Vertical pole\\\hline
\end{comps}

\begin{remglyph}
\hooglig{Intervene} in that \remcom{person's} life with a \remcom{slashing} \remcom{pole}.\\\hline
\end{remglyph}

\begin{prons}
\pronsrtwocone&\textbf{カイ (KAI)}&---\\\hline
\pronsrthreecone&---&---\\\hline
\end{prons}

\begin{remprons}
\hooglig{Intervening} \comon{kites}.\\\hline
\end{remprons}

%{\middel}
%&()\\\hline
\begin{somme}
介入&intervene + enter = \textit{intervention}&カイ{\middel}ニュウ　\mbox{(KAI{\middel}NY\={U})}\\\hline
不介入&not + intervene + enter = \textit{non-intervention}&フ{\middel}カイ{\middel}ニュウ \mbox{(FU{\middel}KAI{\middel}NY\={U})}\\\hline
魚介&fish + intervene = \textit{marine products}&ギョ{\middel}カイ \mbox{(GYO{\middel}KAI)}\\\hline
耳介&ear + intervene = \textit{auricle}&ジ{\middel}カイ \mbox{(JI{\middel}KAI)}\\\hline
米の介入&rice + の + intervention = \textit{intervention of America}&ベイ{\middel}の{\middel}カイ{\middel}ニュウ \mbox{BEI{\middel}no{\middel}KAI{\middel}NY\={U}}\\\hline
\end{somme}
\end{multicols}

\begin{decomp}{6}
彼女&の&プライバシイ&に&介入する&な。\\\hline
かのジョ&の&プライバシイ&に&カイニュウする&な\\\hline
her&の&privacy&に&intervene&な\\\hline
\multicolumn{6}{|r|}{\textit{Don't intrude on her privacy.}}\\\hline
\end{decomp}
\vhrulefill{2pt}
%%--------------------------------------------------------------------
%14-------------------------------------------------------------------
\begin{multicols}{2}
\begin{glyph}{Platform (台)}{platform}\label{glyph:platform}
This 漢字 relates to objects that tend to be \textit{flat and elevated} (think of things such as tables, stands, flatlands, and stages, etc.).\\\hline
It is also used as a \textit{counter for machinery and vehicles}.
\end{glyph}
\gindex{Platform}{台}\strindex{5}{台}

\begin{comps}
\comprtwocone&厶 + 口\\\hline
\comprthreecone&Self/Private + Mouth\\\hline
\end{comps}

\begin{remglyph}
\hooglig{Platform} \remcom{privately} used by \remcom{vampires}.\\\hline
\end{remglyph}

\begin{prons}
\pronsrtwocone&\textbf{ダイ、タイ (DAI, TAI)}&---\\\hline
\rye{3}{ダイ is by far the more common reading.}
\pronsrthreecone&---&---\\\hline
\end{prons}

\begin{remprons}
\hooglig{Platforms} \comon{dyed} \comon{tightly}.\\\hline
\end{remprons}

%{\middel}
%&()\\\hline
\begin{somme}
\rye{3}{Note the mix of 音読み and 訓読み in the first compound below.}
台所&platform + location = \textit{kitchen}&ダイ{\middel}どころ　\mbox{(DAI{\middel}dokoro)}\\\hline
天文台&heaven + literature + platform = \textit{astronomical observatory}&テン{\middel}モン{\middel}ダイ \mbox{(TEN{\middel}MON{\middel}DAI)}\\\hline
\notyet{台風}&platform + wind = \textit{typhoon}&タイ{\middel}フウ \mbox{(TAI{\middel}F\={U})}\\\hline
番台&order (turn) + platform = \textit{watch stand}&バン{\middel}ダイ \mbox{(BAN{\middel}DAI)}\\\hline
\notyet{台頭}&platform + head = \textit{rise of}&タイ{\middel}トウ \mbox{(TAI{\middel}T\={O})}\\\hline
\end{somme}
\end{multicols}

\begin{decomp}{7}
その&天文台&は&良い&位置&に&ある。\\\hline
その&テンモンダイ&は&よい&イチ&に&ある\\\hline
that&observatory&は&good&place&に&to be\\\hline
\multicolumn{7}{|r|}{\textit{That observatory stands in a good location.}}\\\hline
\end{decomp}
\vspace*{\fill}
\pagebreak
%%--------------------------------------------------------------------
%15-------------------------------------------------------------------
\vspace*{\fill}
\begin{multicols}{2}
\begin{glyph}{Hasten (急)}{hasten}\label{glyph:hasten}
Hasten.\\\hline
The sense of things being urgent or sudden.
\end{glyph}
\gindex{Hasten}{急}\strindex{9}{急}

\begin{comps}
\comprtwocone&勹 + 彐 + 心\\\hline
\comprthreecone&Wrap/Embracing + Pig's Head/Snout + Heart\\\hline
\end{comps}

\begin{remglyph}
\hooglig{Hasten} in \remcom{embracing} \remcom{YO} \remcom{heart}.\\\hline
\end{remglyph}

\begin{prons}
\pronsrtwocone&\textbf{キュウ (KY\={U})}&\textbf{いそ (iso)}\\\hline
\pronsrthreecone&---&---\\\hline
\end{prons}

\begin{remprons}
\hooglig{Hasten} to \comkun{isolate} yourself before he \comon{kills you}.\\\hline
\end{remprons}

%{\middel}
%&()\\\hline
\begin{somme}
急ぐ&\textit{hasten; go quickly}&いそ{\middel}ぐ \mbox{(iso{\middel}gu)}\\\hline
急死&hasten + death = \textit{sudden death}&キュウ{\middel}シ \mbox{(KY\={U}{\middel}SHI)}\\\hline
急進&hasten + advance = \textit{rapid progress; radical}&キュウ{\middel}シン \mbox{(KY\={U}{\middel}SHIN)}\\\hline
\notyet{急変}&hasten + alter = \textit{sudden change}&キュウ{\middel}ヘン \mbox{(KY\={U}{\middel}HEN)}\\\hline
\notyet{急降下}&hasten + descend + lower = \textit{nose dive}&キュウ{\middel}コウ{\middel}カ \mbox{(KY\={U}{\middel}K\={O}{\middel}KA)}\\\hline
\notyet{特急}&special + hasten = \textit{limited express}&トッ{\middel}キュウ \mbox{(TOK{\middel}KY\={U})}\\\hline
\end{somme}

\begin{decomp}{2}
急いで、&トム。\\\hline
いそいで&トム\\\hline
hurry&Tom\\\hline
\multicolumn{2}{|r|}{\textit{Hurry up, Tom.}}\\\hline
\end{decomp}

\end{multicols}
\vhrulefill{2pt}
%%--------------------------------------------------------------------
%16-------------------------------------------------------------------
\begin{multicols}{2}
\begin{glyph}{State (州)}{state}\label{glyph:state}
The \textit{political subdivision} of a state or a province.
\end{glyph}
\gindex{State}{州}\strindex{6}{州}

\begin{comps}
\comprtwocone&川 + 丶\\\hline
\comprthreecone&River + Dot\\\hline
\end{comps}

\begin{remglyph}
The \hooglig{state}　that is \remcom{dotted} with \remcom{rivers}.\\\hline
\end{remglyph}

\begin{prons}
\pronsrtwocone&\textbf{シュウ (SH\={U})}&---\\\hline
\pronsrthreecone&---&す (su)\\\hline
\end{prons}

\begin{remprons}
The \hooglig{state} played \ucomkun{sudoku} without their \comon{shoes} on.\\\hline
\end{remprons}

%{\middel}
%&()\\\hline
\begin{somme}
本州&main + state = \textit{Honsh\={u}}&ホン{\middel}シュウ \mbox{(HON{\middel}SH\={U})}\\\hline
九州&nine + state = \textit{Ky\={u}sh\={u}}&キュウ{\middel}シュウ \mbox{(KY\={U}{\middel}SH\={U})}\\\hline
神州&god + state = \textit{land of the gods}&シン{\middel}シュウ \mbox{(SHIN{\middel}SH\={U})}\\\hline
\notyet{州都}&state + metropolis = \textit{state capital}&シュウ{\middel}ト \mbox{(SH\={U}{\middel}TO)}\\\hline
\notyet{州法}&state + law = \textit{state law}&シュウ{\middel}ホウ \mbox{(SH\={U}{\middel}H\={O})}\\\hline
\end{somme}
\end{multicols}

\begin{decomp}{8}
私&は&九州&の&福岡&の&生まれ&です。\\\hline
わたし&は&キュウシュウ&の&フクおか&の&うまれ&です\\\hline
I&は&Ky\={u}sh\={u}&の&Fukuoka&の&birth place&be/is\\\hline
\multicolumn{8}{|r|}{\textit{I'm from Fukuoka in Ky\={u}sh\={u}}.}\\\hline
\end{decomp}
\vspace*{\fill}
\pagebreak
%%--------------------------------------------------------------------
%17-------------------------------------------------------------------
\vspace*{\fill}
\begin{multicols}{2}
\begin{glyph}{Spicy (辛)}{spicy}\label{glyph:spicy}
Spicy.  Pungent.\\\hline
By extension, this 漢字 is also used in some compounds to express the ideas of \textit{hardship} and \textit{trials}.
\end{glyph}
\gindex{Spicy}{辛}\strindex{7}{辛}

\begin{comps}
\comprtwocone&辛\\\hline
\comprthreecone&Spicy/Bitter\\\hline
\end{comps}

\begin{remglyph}
Part of the 漢字 radicals (\#{160}).\\\hline
\end{remglyph}

\begin{prons}
\pronsrtwocone&---&\textbf{から (kara)}\\\hline
\pronsrthreecone&シン (SHIN)&---\\\hline
\end{prons}

\begin{remprons}
\hooglig{Spicy} \comkun{caramel} \ucomon{sheen}.\\\hline
\end{remprons}

%{\middel}
%&()\\\hline
\begin{somme}
辛い&\textit{spicy; hot}&から{\middel}い (kara{\middel}i)\\\hline
辛口&spicy + mouth = \textit{dry (sak\'{e}/wine)}&から{\middel}くち (kara{\middel}kuchi)\\\hline
聞き辛い&ear + spicy = \textit{difficult to hear/ask}&き{\middel}き づら{\middel}い \mbox{(ki{\middel}ki zu{\middel}rai)}\\\hline
分かり辛い&part + spicy = \textit{difficult to understand}&わ{\middel}かり　づら{\middel}い \mbox{(wa{\middel}kari zura{\middel}i)}\\\hline
\notyet{点が辛い}&point + spicy = \textit{severe in marking}＆&テン{\middel}が からい \mbox{(TEN{\middel}ga kara{\middel}i)}\\\hline
\end{somme}

\begin{decomp}{5}
この&水&は&少し&塩辛い。\\\hline
この&みず&は&すこし&しおからい\\\hline
this&water&は&little&salty\\\hline
\multicolumn{5}{|r|}{\textit{This water is a little salty.}}\\\hline
\end{decomp}

\end{multicols}
\vhrulefill{2pt}
%%--------------------------------------------------------------------
%18-------------------------------------------------------------------
\begin{multicols}{2}
\begin{glyph}{Clothes (服)}{clothes}\label{glyph:clothes}
Clothes.\\\hline
One secondary meaning relates to the ideas of \textit{obey} and \textit{submission}.\\\hline
Another secondary meaning has to do with \textit{a dose of something}.
\end{glyph}
\gindex{Clothes}{服}\strindex{8}{服}

\begin{comps}
\comprtwocone&月 + 卩 + 又\\\hline
\comprthreecone&Body/Moon + Seal/Stamp + Again\\\hline
\end{comps}

\begin{remglyph}
\hooglig{Clothed} \remcom{bodies} that are \remcom{sealed} \remcom{again}.\\\hline
\end{remglyph}

\begin{prons}
\pronsrtwocone&\textbf{フク (FUKU)}&---\\\hline
\pronsrthreecone&---&---\\\hline
\end{prons}

\begin{remprons}
\hooglig{Clothes}\ldots \comon{Who cooks} that?\\\hline
\end{remprons}

%{\middel}
%&()\\\hline
\begin{somme}
服&\textit{clothes}&フク \mbox{(FUKU)}\\\hline
一服&one + clothes (dose) = \textit{a dose; a smoke; a rest}&イッ{\middel}プク \mbox{(IP{\middel}PUKU)}\\\hline
学生服&student + clothes = \textit{school uniform}&ガク{\middel}セイ{\middel}フク \mbox{(GAKU{\middel}SEI{\middel}FUKU)}\\\hline
服用&dose + utilize = \textit{taking medicine}&フク{\middel}ヨウ \mbox{(FUKU{\middel}Y\={O})}\\\hline
衣服&clothes + clothes = \textit{clothes}&イ{\middel}フク \mbox{I{\middel}FUKU}\\\hline
\notyet{洋服}&ocean + clothes = \textit{Western-style clothes}&ヨウ{\middel}フク \mbox{(Y\={O}{\middel}FUKU)}\\\hline
\end{somme}

\begin{decomp}{4}
良い&服&を&買え。\\\hline
いい&フク&を&かえ\\\hline
good&clothes&を&buy\\\hline
\multicolumn{4}{|r|}{\textit{Get yourself a decent suit.}}\\\hline
\end{decomp}

\end{multicols}
\vspace*{\fill}
\pagebreak
%%--------------------------------------------------------------------
%19-------------------------------------------------------------------
\vspace*{\fill}
\begin{multicols}{2}
\begin{glyph}{Lucky (吉)}{lucky}\label{glyph:lucky}
Lucky.  Fortunate.
\end{glyph}
\gindex{Lucky}{吉}\strindex{6}{吉}

\begin{comps}
\comprtwocone&士 + 口\\\hline
\comprthreecone&Gentleman + Mouth\\\hline
\end{comps}

\begin{remglyph}
\hooglig{Luckily}, the \remcom{gentleman} wasn't the \remcom{vampire}.\\\hline
\end{remglyph}

\begin{prons}
\pronsrtwocone&\textbf{キチ、キツ (KICHI, KITSU)}&---\\\hline
\pronsrthreecone&---&---\\\hline
\end{prons}

\begin{remprons}
\hooglig{Lucky} \comon{kits} for \comon{key chinking}.\\\hline
\end{remprons}

%{\middel}
%&()\\\hline
\begin{somme}
\rye{3}{This is another of those quirky 漢字 that have multiple common readings and only a few compounds.}
吉日&lucky + sun (day) = \textit{lucky day}&キチ{\middel}ジツ \mbox{(KICHI{\middel}JITSU)}\\\hline
不吉&not + lucky = \textit{unlucky}&フ{\middel}キツ \mbox{(FU{\middel}KITSU)}\\\hline
大吉&big + lucky = \textit{excellent luck}&ダイ{\middel}キチ \mbox{(DAI{\middel}KICHI)}\\\hline
吉凶&lucky + evil = \textit{sunshine and shadow}&キッ{\middel}キョウ \mbox{(KIK{\middel}KY\={O})}\\\hline
\notyet{末吉}&end + lucky = \textit{good luck to come}&すえ{\middel}キチ \mbox{(sue{\middel}KICHI)}\\\hline
\end{somme}
\end{multicols}

\begin{decomp}{11}
十三&は&不吉&な&数&で&ある&と言う&人&が&いる。\\\hline
ジュウサン&は&フキツ&な&かず&で&ある&という&ひと&が&いる\\\hline
13&は&unlucky&な&number&で&some&called&person&が&to be\\\hline
\multicolumn{11}{|r|}{\textit{Some people say thirteen is an unlucky number.}}\\\hline
\end{decomp}
\vhrulefill{2pt}
%%--------------------------------------------------------------------
%20-------------------------------------------------------------------
\begin{multicols}{2}
\begin{glyph}{Utilize (用)}{utilize}\label{glyph:utilize}
Utilize.\\\hline
This important 漢字 also has the sense of a person's \textit{business} or \textit{errands}.\\\hline
It is often translated as \textit{for} when used as the last character in a compound (illustrated by the final example below).
\end{glyph}
\gindex{Utilize}{用}\strindex{5}{用}

\begin{comps}
\comprtwocone&用\\\hline
\comprthreecone&Utilize/Ground\\\hline
\end{comps}

\begin{remglyph}
Part of the 漢字 radicals (\#{101}).\\\hline
\end{remglyph}

\begin{prons}
\pronsrtwocone&\textbf{ヨウ (Y\={O})}&---\\\hline
\pronsrthreecone&---&もち (mochi)\\\hline
\end{prons}

\begin{remprons}
\hooglig{Utilize} \ucomkun{mochi} for \comon{yawning}.\\\hline
{\warn}\textit{Mochi} is a Japanese rice cake.\\\hline
\end{remprons}

%{\middel}
%&\mbox{()}\\\hline
\begin{somme}
用事&utilize + matter = \textit{business; things to be done}&ヨウ{\middel}ジ \mbox{(Y\={O}{\middel}JI)}\\\hline
用心&utilize + heart = \textit{caution; care}&ヨウ{\middel}ジン \mbox{(Y\={O}{\middel}JIN)}\\\hline
男子用&man + child + utilize = \textit{for men; men's}&ダン{\middel}シ{\middel}ヨウ \mbox{(DAN{\middel}SHI{\middel}Y\={O})}\\\hline
用語&utilize + words = \textit{term; terminology}&ヨウ{\middel}ゴ \mbox{(Y\={O}{\middel}GO)}\\\hline
\notyet{利用}&profit + utilize = \textit{use; application}&リ{\middel}ヨウ \mbox{(RI{\middel}Y\={O})}\\\hline
\end{somme}

\begin{decomp}{4}
スリ&に&ご&用心。\\\hline
スリ&に&ご&ヨウジン\\\hline
pickpocket&に&ご&caution\\\hline
\multicolumn{4}{|r|}{\textit{Look out for pickpockets.}}\\\hline
\end{decomp}

\end{multicols}
\vspace*{\fill}
\pagebreak
%%--------------------------------------------------------------------
%21-------------------------------------------------------------------
\vspace*{\fill}
\begin{multicols}{2}
\begin{glyph}{Next (次)}{next}\label{glyph:next}
Next.  Following.
\end{glyph}
\gindex{Next}{次}\strindex{6}{次}

\begin{comps}
\comprtwocone&冫 + 欠\\\hline
\comprthreecone&Ice/Two Dots + Lack\\\hline
\end{comps}

\begin{remglyph}
\hooglig{Following} the \remcom{ice} \remcom{lack}.\\\hline
\end{remglyph}

\begin{prons}
\pronsrtwocone&\textbf{ジ (JI)}&\textbf{つ、つぎ (tsu, tsugi)}\\\hline
\rye{3}{ジ　is the most common of these readings.\\\hline
つ is found only as the verb in the second example below.}
\pronsrthreecone&シ (SHI)&---\\\hline
\rye{3}{This reading is found in only one word, 次第, which we met earlier as the last compound for 第 in Entry \ref{glyph:number} (page \pageref{glyph:number}).}
\end{prons}

\begin{remprons}
\hooglig{Next} to the \ucomon{shield}, was a \comon{jig} of \comkun{stew gifts}.\\\hline
\end{remprons}

%{\middel}
%&\mbox{()}\\\hline
\begin{somme}
次&\textit{next}&つぎ \mbox{(tsugi)}\\\hline
次ぐ&\textit{come after; follow}&つ{\middel}ぐ \mbox{(tsu{\middel}gu)}\\\hline
次々&next + next = \textit{in succession; one by one}&つぎ{\middel}つぎ \mbox{(tsugi{\middel}tsugi)}\\\hline
目次&eye + next = \textit{table of contents}&モク{\middel}ジ \mbox{(MOKU{\middel}JI)}\\\hline
次男&next + man = \textit{second son}&ジ{\middel}ナン \mbox{JI{\middel}NAN}\\\hline
次回&next + rotate = \textit{next time (occasion)}じ&ジ{\middel}カイ \mbox{(JI{\middel}KAI)}\\\hline
次第&next + number = \textit{order; circumstances}&シ{\middel}ダイ \mbox{(SHI{\middel}DAI)}\\\hline
\end{somme}

\begin{decomp}{6}
春&は&冬&の&次&です。\\\hline
はる&は&ふゆ&の&つぎ&です\\\hline
spring&は&winter&の&next&be/is\\\hline
\multicolumn{6}{|r|}{\textit{Spring comes after winter.}}\\\hline
\end{decomp}

\end{multicols}
\vhrulefill{2pt}
%%--------------------------------------------------------------------
%22-------------------------------------------------------------------
\begin{multicols}{2}
\begin{glyph}{Low (低)}{low}\label{glyph:low}
Low.
\end{glyph}
\gindex{Low}{低}\strindex{7}{低}

\begin{comps}
\comprtwocone&亻 + 氏 + 一\\\hline
\comprthreecone&Person + Surname + One\\\hline
\end{comps}

\begin{remglyph}
\hooglig{Low} \remcom{people} belong to \remcom{one} \remcom{clan}.\\\hline
\end{remglyph}

\begin{prons}
\pronsrtwocone&\textbf{テイ (TEI)}&\textbf{ひく (hiku)}\\\hline
\pronsrthreecone&---&---\\\hline
\end{prons}

\begin{remprons}
\hooglig{Lower} your \comkun{hiccups} with this \comon{tape}.\\\hline
\end{remprons}

%{\middel}
%&\mbox{()}\\\hline
\begin{somme}
低い&\textit{low}&ひく{\middel}い \mbox{(hiku{\middel}i)}\\\hline
低下&low + lower = \textit{lowering; decline}&テイ{\middel}カ \mbox{(TEI{\middel}KA)}\\\hline
最低&most + low = \textit{lowest; worst}&サイ{\middel}テイ \mbox{(SAI{\middel}TEI)}\\\hline
低音&low + sound = \textit{low tone; bass; low voice}&テイ{\middel}オン \mbox{(TEI{\middel}ON)}\\\hline
\notyet{低温}&low + warm = \textit{low temperature}&テイ{\middel}オン \mbox{(TEI{\middel}ON)}\\\hline
\notyet{低速}&low + fast = \textit{low speed}&テイ{\middel}ソク \mbox{(TEI{\middel}SOKU)}\\\hline
\end{somme}

\begin{decomp}{6}
私&は&血圧&が&低い&です。\\\hline
わたし&は&ケツアツ&が&ひくい&です\\\hline
I&は&blood pressure&が&low&be/is\\\hline
\multicolumn{6}{|r|}{\textit{I have low blood pressure.}}\\\hline
\end{decomp}

\end{multicols}
\vspace*{\fill}
\pagebreak
%%--------------------------------------------------------------------
%23-------------------------------------------------------------------
\vspace*{\fill}
\begin{multicols}{2}
\begin{glyph}{Filial Piety (孝)}{filial_piety}\label{glyph:filial_piety}
Filial piety.  Obedience to one's parents.
\end{glyph}
\gindex{Filial Piety}{孝}\strindex{7}{孝}

\begin{comps}
\comprtwocone&老 + 子\\\hline
\comprthreecone&Old + Child\\\hline
\end{comps}

\begin{remglyph}
\hooglig{Filial piety} the \remcom{old} requires from \remcom{children}.\\\hline
\end{remglyph}

\begin{prons}
\pronsrtwocone&\textbf{コウ (K\={O})}&---\\\hline
\pronsrthreecone&---&---\\\hline
\end{prons}

\begin{remprons}
\hooglig{Filial piety} is a noble \comon{cause}.\\\hline
\end{remprons}

%{\middel}
%&\mbox{()}\\\hline
\begin{somme}
\rye{3}{Note the mix of 音読み and 訓読み in two sample compounds.}
親孝行&parent + filial piety + go = \textit{filial piety}&おや{\middel}コウ{\middel}コウ \mbox{(oya{\middel}K\={O}{\middel}K\={O})}\\\hline
親不孝&parent + not + filial piety = \textit{lack of filial piety}&おや{\middel}フ{\middel}コウ \mbox{(oya{\middel}FU{\middel}K\={O})}\\\hline
至孝&reach + filial piety = \textit{supreme filial piety}&シ{\middel}コウ \mbox{(SHI{\middel}K\={O})}\\\hline
孝弟&filial piety + younger brother = \textit{brotherly love}&コウ{\middel}テイ \mbox{(K\={O}{\middel}TEI)}\\\hline
孝心&filial piety + heart = \textit{filial devotion}&コウ{\middel}シン \mbox{(K\={O}{\middel}SHIN)}\\\hline
\end{somme}
\end{multicols}

\begin{decomp}{5}
親孝したい&時に&親&は&無し。\\\hline
おやコウしたい&ときに&おや&は&ない\\\hline
filial piety&incidentally&parents&は&non existent\\\hline
\multicolumn{5}{|r|}{\textit{When one would be filial, one's parents are gone.}}\\\hline
\end{decomp}
\vhrulefill{2pt}
%%--------------------------------------------------------------------
%24-------------------------------------------------------------------
\begin{multicols}{2}
\begin{glyph}{Tea (茶)}{tea}\label{glyph:tea}
Tea.
\end{glyph}
\gindex{Tea}{茶}\strindex{9}{茶}

\begin{comps}
\comprtwocone&艹 + 人 + ホ\\\hline
\comprthreecone&Grass/Herb/Plant + Person + HO (tree)\\\hline
\end{comps}

\begin{remglyph}
\hooglig{Tea} \remcom{herbs} \remcom{people} got from \remcom{trees}.\\\hline
\end{remglyph}

\begin{prons}
\pronsrtwocone&\textbf{チャ (CHA)}&---\\\hline
\pronsrthreecone&サ (SA)&---\\\hline
\rye{3}{Though not common, this reading does show up in the Japanese word for \textit{tea ceremony}, 茶道, サ{\middel}ドウ/SA{\middel}D\={O}.}
\end{prons}

\begin{remprons}
\hooglig{Tea} \comon{chums} \ucomon{subjecting} themselves to tea experiments.\\\hline
\end{remprons}

%{\middel}
%&\mbox{()}\\\hline
\begin{somme}
茶&\textit{tea}&チャ　\mbox{(CHA)}\\\hline
茶会&tea + meet = \textit{tea party; get together}&チャ{\middel}カイ/サ{\middel}カイ \mbox{(CHA{\middel}KAI/SA{\middel}KAI)}\\\hline
無茶&nothingness + tea = \textit{absurd}&ム{\middel}チャ/ムッ{\middel}チャ \mbox{(MU(C){\middel}CHA)}\\\hline
\notyet{茶屋}&tea + roof = \textit{tea house; tea dealer}&チャ{\middel}や \mbox{(CHA{\middel}ya)}\\\hline
\notyet{茶色}&tea + colour = \textit{brown}&チャ{\middel}いろ \mbox{(CHA{\middel}iro)}\\\hline
\notyet{茶室}&tea + room = \textit{tea room}&チャ{\middel}シツ \mbox{(CHA{\middel}SHITSU)}\\\hline
\end{somme}

\begin{decomp}{4}
あれ&は&茶番&だ。\\\hline
あれ&は&チャバン&だ\\\hline
that&は&farce&だ\\\hline
\multicolumn{4}{|r|}{\textit{That's a farce.}}\\\hline
\end{decomp}

\end{multicols}
\vspace*{\fill}
\pagebreak
%%--------------------------------------------------------------------
%25-------------------------------------------------------------------
\vspace*{\fill}
\begin{multicols}{2}
\begin{glyph}{Profit (利)}{profit}\label{glyph:profit}
Profit.  Advantage.
\end{glyph}
\gindex{Profit}{利}\strindex{7}{利}

\begin{comps}
\comprtwocone&禾 + 刂\\\hline
\comprthreecone&Grain Ear + Knife\\\hline
\end{comps}

\begin{remglyph}
\hooglig{Profit} equals the \remcom{grain} \remcom{knifed} off.\\\hline
\end{remglyph}

\begin{prons}
\pronsrtwocone&\textbf{リ (RI)}&\textbf{き (ki)}\\\hline
\pronsrthreecone&---&---\\\hline
\end{prons}

\begin{remprons}
\hooglig{Profits} \comkun{keep} \comon{leaping}.\\\hline
\end{remprons}

%{\middel}
%&\mbox{()}\\\hline
\begin{somme}
利く&\textit{take effect; work}&き{\middel}く \mbox{(ki{\middel}ku)}\\\hline
利子&profit + child = \textit{profit (on money)}&リ{\middel}シ \mbox{(RI{\middel}SHI)}\\\hline
利口&profit + mouth = \textit{smart; clever}&リ{\middel}コウ \mbox{(RI{\middel}K\={O})}\\\hline
利用&profit + utilize = \textit{use; application}&リ{\middel}ヨウ \mbox{(RI{\middel}Y\={O})}\\\hline
有利&have + profit = \textit{advantageous; profitable}&ユウ{\middel}リ \mbox{(Y\={U}{\middel}RI)}\\\hline
利点&profit + point = \textit{point in favour}&リ{\middel}テン \mbox{(RI{\middel}TEN)}\\\hline
利き手&profit + hand = \textit{one's dominant hand}&き{\middel}きて \mbox{(ki{\middel}kite)}\\\hline
\end{somme}

\begin{decomp}{4}
ジョン&は&利口&だ。\\\hline
ジョン&は&リコウ&だ\\\hline
John&は&clever&だ\\\hline
\multicolumn{4}{|r|}{\textit{John is clever.}}\\\hline
\end{decomp}

\end{multicols}
\vhrulefill{2pt}
%%--------------------------------------------------------------------
%26-------------------------------------------------------------------
\begin{multicols}{2}
\begin{glyph}{Not Yet (未)}{not_yet}\label{glyph:not_yet}
Not yet.\\\hline
Invariably used as the first glyph in a compound to indicate something \textit{incomplete} or \textit{unfinished}.
\end{glyph}
\gindex{Not Yet}{未}\strindex{5}{未}

\begin{comps}
\comprtwocone&一 + 木\\\hline
\comprthreecone&One + Tree\\\hline
\end{comps}

\begin{remglyph}
\hooglig{Incomplete} and \remcom{one} lonely \remcom{tree}.\\\hline
\end{remglyph}

\begin{prons}
\pronsrtwocone&\textbf{ミ (MI)}&---\\\hline
\pronsrthreecone&---&---\\\hline
\end{prons}

\begin{remprons}
\hooglig{Not yet} \comon{mitigated}.\\\hline
\end{remprons}

%{\middel}
%&\mbox{()}\\\hline
\begin{somme}
未来　&not yet + come = \textit{future; future tense}&ミ{\middel}ライ \mbox{(MI{\middel}RAI)}\\\hline 
未定&not yet + fixed = \textit{undecided; pending}&ミ{\middel}テイ \mbox{(MI{\middel}TEI)}\\\hline
未明&not yet + bright = \textit{early dawn}&ミ{\middel}メイ \mbox{(MI{\middel}MEI)}\\\hline
未亡人&not yet + perish + person = \textit{widow}&ミ{\middel}ボウ{\middel}ジン \mbox{(MI{\middel}B\={O}{\middel}JIN)}\\\hline
\notyet{未完}&not yet + perfect = \textit{incomplete; unfinished}&ミ{\middel}カン \mbox{(MI{\middel}KAN)}\\\hline
\notyet{未婚}&not yet + marriage = \textit{unmarried}&ミ{\middel}コン \mbox{(MI{\middel}KON)}\\\hline
\notyet{未開発}&not yet + open + discharge = \textit{undeveloped}&ミ{\middel}カイ{\middel}ハツ \mbox{(MI{\middel}KAI{\middel}HATSU)}\\\hline
\end{somme}

\begin{decomp}{7}
君&に&は&明るい&未来&が&ある。\\\hline
きみ&に&は&あかるい&ミライ&が&ある\\\hline
you&に&は&bright&future&が&to be\\\hline
\multicolumn{7}{|r|}{\textit{You have a bright future.}}\\\hline
\end{decomp}

\end{multicols}
\vspace*{\fill}
\pagebreak
%%--------------------------------------------------------------------
%27-------------------------------------------------------------------
\vspace*{\fill}
\begin{multicols}{2}
\begin{glyph}{End (末)}{end}\label{glyph:end}
The \textit{end of things}, in space as well as in time.
\end{glyph}
\gindex{End}{末}\strindex{5}{末}

\begin{comps}
\comprtwocone&一 + 木\\\hline
\comprthreecone&One + Tree\\\hline
\end{comps}

\begin{remglyph}
\hooglig{End} that \remcom{one} \remcom{tree}.\\\hline
\end{remglyph}

\begin{prons}
\pronsrtwocone&\textbf{マツ (MATSU)}&---\\\hline
\pronsrthreecone&バツ (BATSU)&すえ (sue)\\\hline
\end{prons}

\begin{remprons}
\hooglig{End} the \ucomkun{suede} \comon{mutt} \ucomon{butts}.\\\hline
{\warn}\textit{Suede} is a type of leather.\\\hline
\end{remprons}

%{\middel}
%&\mbox{()}\\\hline
\begin{somme}
週末　&week + end = \textit{weekend}&シュウ{\middel}マツ \mbox{(SH\={U}{\middel}MATSU)}\\\hline
年末&year + end = \textit{end of the year}&ネン{\middel}マツ \mbox{(NEN{\middel}MATSU)}\\\hline
月末&moon + end = \textit{end of the month}&ゲツ{\middel}マツ \mbox{(GETSU{\middel}MATSU)}\\\hline
末日&end + day = \textit{last day (of a month)}&マツ{\middel}ジツ \mbox{(MATSU{\middel}JITSU)}\\\hline
\notyet{末広}&end + wide = \textit{spread out like a fan; becoming prosperous; folding fan}&すえ{\middel}ひろ \mbox{(sue{\middel}hiro)}\\\hline
\notyet{末期}&end + period = \textit{last stage}&マッ{\middel}キ \mbox{(MAK{\middel}KI)}\\\hline
\end{somme}

\begin{decomp}{3}
良い&週末&を。\\\hline
よい&シュウマツ&を\\\hline
good&weekend&を\\\hline
\multicolumn{3}{|r|}{\textit{Have a nice weekend!}}\\\hline
\end{decomp}

\end{multicols}
\vhrulefill{2pt}
%%--------------------------------------------------------------------
%28-------------------------------------------------------------------
\begin{multicols}{2}
\begin{glyph}{Dipper (斗)}{dipper}\label{glyph:dipper}
Dipper.\\\hline
A kind of volume measure.
\end{glyph}
\gindex{Dipper}{斗}\strindex{4}{斗}

\begin{comps}
\comprtwocone&斗\\\hline
\comprthreecone&Dipper/A kind of volume measure\\\hline
\end{comps}

\begin{remglyph}
Part of the 漢字 radicals (\#{68}).\\\hline
\end{remglyph}

\begin{prons}
\pronsrtwocone&\textbf{ト (TO)}&---\\\hline
\pronsrthreecone&---&---\\\hline
\end{prons}

\begin{remprons}
\hooglig{Dipper} at the \comon{top} of the sky.\\\hline
\end{remprons}

%{\middel}
%&\mbox{()}\\\hline
\begin{somme}
北斗七星&North + Dipper + seven + star = \textit{the Big Dipper}&ホク{\middel}ト{\middel}シチ{\middel}セイ (HOKU{\middel}TO{\middel}SHICHI{\middel}SEI)\\\hline
斗酒&volume measure + sak\'{e} = \textit{kegs of sak\'{e}; lots of sak\'{e}}&ト{\middel}シュ \mbox{(TO{\middel}SHU)}\\\hline
\end{somme}
\end{multicols}

\begin{decomp}{5}
北斗七星&は&簡単&に&見つからん。\\\hline
ホクトシチセイ&は&カンタン&に&みつから\\\hline
Big dipper&は&simple&に&to be found\\\hline
\multicolumn{5}{|r|}{\textit{You can find the Big Dipper easily.}}\\\hline
\end{decomp}
\vspace*{\fill}
\pagebreak
%%--------------------------------------------------------------------
%29-------------------------------------------------------------------
\vspace*{\fill}
\begin{multicols}{2}
\begin{glyph}{Play (遊)}{play}\label{glyph:play}
Play.  Amusement.
\end{glyph}
\gindex{Play}{遊}\strindex{12}{遊}

\begin{comps}
\comprtwocone&辶 + 方 + 子 + ⺅\\\hline
\comprthreecone&Road/Walk + Direction + Child + Person\\\hline
\end{comps}

\begin{remglyph}
\hooglig{Play} in a \remcom{way} with that \remcom{person's} \remcom{child} on the \remcom{road}.\\\hline
\end{remglyph}

\begin{prons}
\pronsrtwocone&\textbf{ユウ (Y\={U})}&\textbf{あそ (aso)}\\\hline
\pronsrthreecone&ユ (YU)&---\\\hline
\end{prons}

\begin{remprons}
\hooglig{Play} \ucomon{youthfully} on \comon{YouTube} and be \comkun{asocial}.\\\hline
\end{remprons}

%{\middel}
%&\mbox{()}\\\hline
\begin{somme}
遊ぶ&\textit{to play}&あそ{\middel}ぶ \mbox{(aso{\middel}bu)}\\\hline
遊歩道&play + walk + road = \textit{promenade}&ユウ{\middel}ホ{\middel}ドウ \mbox{(Y\={U}{\middel}HO{\middel}D\={O})}\\\hline
遊女&play + woman = \textit{prostitute; harlot}&ユウ{\middel}ジョ \mbox{(Y\={U}{\middel}JO)}\\\hline
周遊&around + play = \textit{excursion ticket; tour; round trip}&シュウ{\middel}ユウ \mbox{(SH\={U}{\middel}Y\={U})}\\\hline
交遊&mix + play = \textit{friend; companion}&コウ{\middel}ユウ \mbox{(K\={O}{\middel}Y\={U})}\\\hline
\notyet{遊泳}&play + swim = \textit{swimming}&ユウ{\middel}エイ \mbox{(Y\={U}{\middel}EI)}\\\hline
\notyet{遊園地}&play + garden + ground = \textit{amusement park; playground}&ユウ{\middel}エン{\middel}チ \mbox{(Y\={U}{\middel}EN{\middel}CHI)}\\\hline
\end{somme}

\begin{decomp}{5}
何&を&して&遊ぼう&か。\\\hline
なに&を&して&あそぼう&か\\\hline
what&を&to do&play&か\\\hline
\multicolumn{5}{|r|}{\textit{What shall we play?}}\\\hline
\end{decomp}

\end{multicols}
\vhrulefill{2pt}
%%--------------------------------------------------------------------
%30-------------------------------------------------------------------
\begin{multicols}{2}
\begin{glyph}{Bean (豆)}{bean}\label{glyph:bean}
Bean.
\end{glyph}
\gindex{Bean}{豆}\strindex{7}{豆}

\begin{comps}
\comprtwocone&豆\\\hline
\comprthreecone&Bean\\\hline
\end{comps}

\begin{remglyph}
Part of the 漢字 radicals (\#{151}).\\\hline
\end{remglyph}

\begin{prons}
\pronsrtwocone&---&\textbf{まめ (mame)}\\\hline
\pronsrthreecone&ズ、トウ (ZU, T\={O})&---\\\hline
\rye{3}{ズ　appears only in the word for \textit{soybean} 大豆 (DAI{\middel}ZU/ダイ{\middel}ズ) made from the 漢字 for `large' and `bean'.\\\hline
トウ is the pronunciation for when this 漢字 shows up in the famous Japanese word \textit{t\={o}fu} 豆腐\notcov (トウ{\middel}フ/T\={O}{\middel}FU), formed from the 漢字 `bean' and `rot'.}
\end{prons}

\begin{remprons}
\hooglig{Beans} \comkun{maimed} by \ucomon{zucchini} in a \ucomon{torturous} way.\\\hline
\end{remprons}

%{\middel}
%&\mbox{()}\\\hline
\begin{somme}
豆&\textit{bean}&まめ \mbox{(mame)}\\\hline
豆腐{\notcov}&bean + rot = \textit{t\={o}fu}&トウ{\middel}フ \mbox{(T\={O}{\middel}FU)}\\\hline
大豆&large + bean = \textit{soybean}&ダイ{\middel}ズ \mbox{(DAI{\middel}ZU)}\\\hline
豆本&bean + main = \textit{miniature book}&まめ{\middel}ホン \mbox{(mame{\middel}HON)}\\\hline
南京豆&south + capital + bean = \textit{peanuts}&ナン{\middel}キン{\middel}まめ \mbox{(NAN{\middel}KIN{\middel}mame)}\\\hline
年の豆&year + bean = \textit{beans of the bean-scattering ceremony}&とし{\middel}の{\middel}まめ \mbox{(toshi{\middel}no{\middel}mame)}\\\hline
\end{somme}

\begin{decomp}{7}
エミリー&は&夕食&に&豆腐&を&食べる。\\\hline
エミリー&は&ユウショク&に&トウフ&を&たべる\\\hline
Emily&は&dinner&に&t\={o}fu&を&eat\\\hline
\multicolumn{7}{|r|}{\textit{Emily ate t\={o}fu at dinner.}}\\\hline
\end{decomp}

\end{multicols}
\vspace*{\fill}
\pagebreak
%%--------------------------------------------------------------------
%---------------------------------------------------------------------
